\chapterhead{3}{Guess First}{Estimating before you measure, then correcting}\label{ch3}

Before you measure anything, you can guess. A guess is not a wild shot. It is a quick, rough answer that tells you what size of answer to expect.

Why bother? Suppose you measure a pencil and your ruler says 1.7~m. If you had guessed first, you would know at once that something is off, because a pencil is not as long as a person's arm. A guess is a smoke alarm for silly slips.

\begin{tryit}{Guess, then measure}
Pick three things in the room. Write a guess for each length in centimetres. Then measure and see how far off you were.
\par\medskip
\renewcommand{\arraystretch}{1.8}
\begin{tabularx}{\linewidth}{|>{\sffamily\small}p{2.2cm}|X|X|X|}
\hline
\textbf{Thing} & \textbf{My guess (cm)} & \textbf{Measured (cm)} & \textbf{Off by (guess $-$ measured)}\\ \hline
 & & & \\ \hline
 & & & \\ \hline
 & & & \\ \hline
\end{tabularx}
\end{tryit}

A negative ``off by'' means you guessed too small. A positive one means you guessed too big. Over time, most people's guesses get better, because they learn what 10~cm, 30~cm or 1~m looks like.

\subhead{Benchmarks: rulers you carry in your head}
A \textbf{benchmark} is a thing whose size you know well. When you have no ruler, you can compare an unknown length to a benchmark.

\begin{center}
\sffamily\small\renewcommand{\arraystretch}{1.5}
\begin{tabular}{@{}ll@{}}
\toprule
\textbf{Benchmark} & \textbf{About how big}\\
\midrule
A bank card or bus pass card & 8.6~cm by 5.4~cm\\
A sheet of letter paper & 21.6~cm by 27.9~cm (A4 paper: 21.0~cm by 29.7~cm)\\
A new, unsharpened pencil & about 19~cm\\
A room door & about 2~m high\\
\bottomrule
\end{tabular}
\end{center}
Real things vary a little, so check yours with a ruler once and remember the number.

\begin{tryit}{Estimate with a sheet of paper}
Lay sheets of paper end to end along a desk, long side along the desk. Count how many sheets and part-sheets it takes. Multiply by 27.9~cm if you use letter paper, or by 29.7~cm if you use A4. Write the estimate, then check with your ruler.\\[4pt]
Sheets: \wline[1.6cm] \quad Estimate: \wline[1.6cm]~cm \quad Ruler: \wline[1.6cm]~cm
\end{tryit}

\subhead{Beans in a jar}
Now for a bigger guess. How many beans are in a jar? You could count them all, but if the jar is big that takes ages. Instead we can \textbf{estimate}: count a small scoop, and use the scoop to work out the whole jar.

\begin{tryit}{Estimate the beans}
\textbf{You need:} a cup of dried beans, a small cup, a jar or tub, a spoon or ruler to level the cup.
\begin{enumerate}[leftmargin=1.6em]
\item Fill the small cup with beans and level the top. Tip them out and \textbf{count} them. That is one scoop.
\item Pour beans from the small cup into the jar, one level cup at a time, until the jar is full. Count the cups.
\item Multiply: beans in one scoop $\times$ number of cups $=$ \textbf{estimate}.
\item Now take two more scoops and count each. Average the three scoops (add them, divide by 3).
\item Multiply the \textbf{average} by the number of cups. This is your corrected estimate.
\item Only now, count every bean in the jar and compare.
\end{enumerate}
\end{tryit}
\adult{Beans can be a choking hazard for little children and pets. Count them at a table and put them away afterwards. Do not eat them.}

\subhead{Practice numbers}
\begin{worked}{Made-up numbers for practice, not from a real jar}
\textbf{One scoop} had 24 beans. The jar took \textbf{9} level cups.\\
First estimate: $24\times9=216$.\\
Two more scoops: 28 and 23. Three scoops: 24, 28, 23. Total $24+28+23=75$. \textbf{Average:} $75\div3=25$.\\
Corrected estimate: $25\times9=225$. The correction is $+9$ beans.\\
Real count: 231. The corrected estimate was off by $231-225=6$.
\end{worked}

\begin{center}
\begin{tikzpicture}[scale=0.92]
  \node[cir] (s1) at (0,1.2) {24};
  \node[cir] (s2) at (0,0) {28};
  \node[cir] (s3) at (0,-1.2) {23};
  \node[bigcir,minimum size=3.2cm] (big) at (0,0) {};
  \node[cir] at (0,1.0) {24}; \node[cir] at (0,0) {28}; \node[cir] at (0,-1.0) {23};
  \node[valcir] (avg) at (3.8,0) {25};
  \node[note] at (3.8,-0.75) {average scoop};
  \draw[popA] (big) -- node[above,note] {pop} (avg);
  \node[note,anchor=east,align=right,text width=2.6cm] at (-1.9,0) {\textbf{kept:} the typical scoop\\\textbf{lost:} which scoop was which};
  \node[cir] (old) at (7.0,0.9) {216};
  \node[wrapcir,minimum size=2.35cm] (w) at (7.0,0.9) {};
  \node[valcir] (new) at (7.0,-1.4) {225};
  \draw[popA] (w) -- node[right,note] {$\times$ using the average} (new);
  \node[note] at (7.0,2.35) {old estimate, wrapped};
  \node[note,font=\sffamily\scriptsize] at (7.0,0.2) {$+9$};
\end{tikzpicture}
\end{center}

\trust{We took three scoops instead of one, and we compared the corrected estimate, 225, with a full count, 231.}
{The average of three scoops is steadier than one scoop, and the final answer was only 6 away from the real count, which is under 3 beans in every hundred.}
{That a different jar, a different kind of bean or a heaped cup would give the same result. One careful jar shows that the method \emph{can} work, not that it always does.}

\subhead{The ladder from seen to concluded}
Scientists keep four kinds of statement apart. We will use this ladder in every science book.

\begin{center}\ladder[0.8]\end{center}

In the bean jar, ``there were 24 beans in my first scoop'' is a \textbf{count}. ``There are about 225 beans in the jar'' is an \textbf{estimate}, because we counted a few and reasoned about all. Notice the word ``about''. We do not say there are exactly 225.

\begin{keepbox}
We used three scoops to estimate the whole jar. \textbf{We kept} \wline[3cm], \textbf{and we set aside} \wline[3.5cm].
\end{keepbox}

\begin{multix}[wow]\multixbot[wow]\tcblower
Three scoops gave 24, 28 and 23, and the jar holds 9 cups. So $24+28+23=75$ beans in the jar! That was easy. Why do you think the real count was so much bigger?
\end{multix}

\subhead{No peeking}
Do not count the whole jar first and then ``estimate''. An estimate made after peeking at the answer proves nothing. Write the estimate down \emph{before} the full count, so the correction is honest.

\subhead{Pacing it out}
You do not need a tape measure to estimate a long distance. You can walk it. First find out how long \emph{your} walking step is, then count steps.

\begin{tryit}{Pace a hallway}
\textbf{You need:} a long tape measure or a long piece of string and a ruler, plus a hallway or a path.
\begin{enumerate}[leftmargin=1.6em]
\item Walk 10 normal steps in a straight line. Measure the distance. Divide by 10 to get the length of one pace.
\item Do it three times and average the three distances. Your pace length is that average divided by 10.
\item Walk the hallway and count your paces. Multiply the paces by your pace length.
\item Measure the hallway with the tape and compare.
\end{enumerate}
\end{tryit}

\begin{worked}{Practice numbers, not from a real hallway}
\textbf{First try.} 10 paces measured 7.2~m, so one pace is $0.72$~m. The hallway took 31 paces. $31\times0.72=22.32$, so \textbf{about 22~m}.\\
\textbf{Corrected.} Two more tries: 7.7~m and 7.4~m. Total $7.2+7.7+7.4=22.3$, and $22.3\div3\approx7.43$~m for 10 paces. One pace is $0.743$~m. $31\times0.743\approx23.0$, so \textbf{about 23~m}.\\
\textbf{Tape measure.} 22.9~m. The first estimate was off by about 0.6~m. The corrected one was off by about 0.1~m.
\end{worked}

\begin{keepbox}
When we walked a hallway to measure it, \textbf{we kept} \wline[3.2cm] \textbf{and we set aside} \wline[3.2cm].
\end{keepbox}

\clearpage

\begin{puzzle}{Which guess was closer?}
A jar really had 231 beans. Three people guessed.
\begin{center}\sans
Priya: 190 \qquad Sam: 260 \qquad Tomas: 225
\end{center}
For each, work out how far off the guess was.\\
Priya \wline[1.5cm] \quad Sam \wline[1.5cm] \quad Tomas \wline[1.5cm]\\
Whose guess was closest? \wline[3cm]
\end{puzzle}

\begin{puzzle}{Fill the jar}
Four scoops of lentils had 18, 20, 19 and 19 lentils. The jar holds 12 scoops.\\
(a) Total of the four scoops: \wline[1.8cm]\\
(b) Average scoop: \wline[1.8cm]\\
(c) Estimate for the jar: \wline[1.8cm]
\end{puzzle}

\begin{puzzle}{Scoop detective}
Four scoops of beans were counted: \textbf{24, 28, 23 and 52}.\\
One count looks strange. Which one? \wline[2cm]\\
Write one way it might have happened, and what you would do about it.\\
\rule{\linewidth}{0.5pt}
\end{puzzle}

\begin{puzzle}{Smoke alarm}
Each measurement has a slip. Use a quick guess to spot it.
\begin{center}\sans\small\renewcommand{\arraystretch}{1.7}
\begin{tabular}{@{}ll@{}}
(a) ``My pencil is 175 cm long.'' & Slip? \wline[3.5cm]\\
(b) ``This classroom door is 2 mm tall.'' & Slip? \wline[3.5cm]\\
(c) ``My shoe is 25 m long.'' & Slip? \wline[3.5cm]\\
\end{tabular}
\end{center}
\end{puzzle}

\clearpage

\begin{puzzle}{Dots in a field}
There are far too many dots to count one by one. The field is split into 8 equal squares, A to H. A to D are along the top and E to H along the bottom.

\begin{center}
\begin{tikzpicture}[x=1.14cm,y=1.14cm]
  \draw[line width=1.4pt,fill=white] (0,0) rectangle (12,6);
  \foreach \x in {3,6,9}{\draw[line width=0.9pt,ink!60] (\x,0) -- (\x,6);}
  \draw[line width=0.9pt,ink!60] (0,3) -- (12,3);
  \node[font=\sffamily\bfseries\small,text=sky!70!black] at (1.5,6.3) {A};
  \node[font=\sffamily\bfseries\small,text=sky!70!black] at (4.5,6.3) {B};
  \node[font=\sffamily\bfseries\small,text=sky!70!black] at (7.5,6.3) {C};
  \node[font=\sffamily\bfseries\small,text=sky!70!black] at (10.5,6.3) {D};
  \node[font=\sffamily\bfseries\small,text=sky!70!black] at (1.5,-0.3) {E};
  \node[font=\sffamily\bfseries\small,text=sky!70!black] at (4.5,-0.3) {F};
  \node[font=\sffamily\bfseries\small,text=sky!70!black] at (7.5,-0.3) {G};
  \node[font=\sffamily\bfseries\small,text=sky!70!black] at (10.5,-0.3) {H};
  \fill[ink] (5.448,3.335) circle (1.1mm);
  \fill[ink] (10.921,2.808) circle (1.1mm);
  \fill[ink] (7.507,4.641) circle (1.1mm);
  \fill[ink] (1.292,1.899) circle (1.1mm);
  \fill[ink] (1.252,4.734) circle (1.1mm);
  \fill[ink] (8.244,0.435) circle (1.1mm);
  \fill[ink] (11.593,5.603) circle (1.1mm);
  \fill[ink] (7.786,3.647) circle (1.1mm);
  \fill[ink] (2.027,0.284) circle (1.1mm);
  \fill[ink] (6.329,0.533) circle (1.1mm);
  \fill[ink] (2.406,1.555) circle (1.1mm);
  \fill[ink] (0.549,2.798) circle (1.1mm);
  \fill[ink] (5.310,4.918) circle (1.1mm);
  \fill[ink] (6.222,3.786) circle (1.1mm);
  \fill[ink] (5.505,1.758) circle (1.1mm);
  \fill[ink] (11.773,5.776) circle (1.1mm);
  \fill[ink] (9.947,4.164) circle (1.1mm);
  \fill[ink] (3.857,1.486) circle (1.1mm);
  \fill[ink] (3.553,0.593) circle (1.1mm);
  \fill[ink] (10.020,2.364) circle (1.1mm);
  \fill[ink] (11.313,4.945) circle (1.1mm);
  \fill[ink] (0.206,1.374) circle (1.1mm);
  \fill[ink] (11.572,2.426) circle (1.1mm);
  \fill[ink] (1.047,3.725) circle (1.1mm);
  \fill[ink] (9.231,1.711) circle (1.1mm);
  \fill[ink] (11.383,4.445) circle (1.1mm);
  \fill[ink] (1.569,1.580) circle (1.1mm);
  \fill[ink] (1.372,0.535) circle (1.1mm);
  \fill[ink] (9.445,1.195) circle (1.1mm);
  \fill[ink] (6.688,2.706) circle (1.1mm);
  \fill[ink] (2.412,4.299) circle (1.1mm);
  \fill[ink] (1.719,3.805) circle (1.1mm);
  \fill[ink] (1.551,2.556) circle (1.1mm);
  \fill[ink] (2.669,1.711) circle (1.1mm);
  \fill[ink] (11.463,4.699) circle (1.1mm);
  \fill[ink] (3.728,5.155) circle (1.1mm);
  \fill[ink] (2.644,2.408) circle (1.1mm);
  \fill[ink] (10.111,3.794) circle (1.1mm);
  \fill[ink] (1.364,5.740) circle (1.1mm);
  \fill[ink] (9.163,2.042) circle (1.1mm);
  \fill[ink] (1.245,3.463) circle (1.1mm);
  \fill[ink] (4.512,2.738) circle (1.1mm);
  \fill[ink] (6.865,5.053) circle (1.1mm);
  \fill[ink] (2.321,1.063) circle (1.1mm);
  \fill[ink] (10.738,4.780) circle (1.1mm);
  \fill[ink] (8.777,5.466) circle (1.1mm);
  \fill[ink] (2.480,5.521) circle (1.1mm);
  \fill[ink] (10.433,3.580) circle (1.1mm);
  \fill[ink] (5.089,0.782) circle (1.1mm);
  \fill[ink] (0.649,5.591) circle (1.1mm);
  \fill[ink] (3.181,4.813) circle (1.1mm);
  \fill[ink] (7.119,1.843) circle (1.1mm);
  \fill[ink] (0.998,1.479) circle (1.1mm);
  \fill[ink] (10.841,1.342) circle (1.1mm);
  \fill[ink] (0.392,1.707) circle (1.1mm);
  \fill[ink] (5.370,0.539) circle (1.1mm);
  \fill[ink] (2.245,2.265) circle (1.1mm);
  \fill[ink] (6.837,0.937) circle (1.1mm);
  \fill[ink] (4.401,5.189) circle (1.1mm);
  \fill[ink] (11.574,3.879) circle (1.1mm);
  \fill[ink] (8.218,3.473) circle (1.1mm);
  \fill[ink] (0.408,5.297) circle (1.1mm);
  \fill[ink] (8.331,5.592) circle (1.1mm);
  \fill[ink] (0.447,3.763) circle (1.1mm);
  \fill[ink] (5.794,4.291) circle (1.1mm);
  \fill[ink] (3.899,5.796) circle (1.1mm);
  \fill[ink] (1.073,3.258) circle (1.1mm);
  \fill[ink] (8.750,4.141) circle (1.1mm);
  \fill[ink] (9.402,5.324) circle (1.1mm);
  \fill[ink] (4.281,4.037) circle (1.1mm);
  \fill[ink] (10.650,5.078) circle (1.1mm);
  \fill[ink] (5.039,4.627) circle (1.1mm);
  \fill[ink] (10.216,3.408) circle (1.1mm);
  \fill[ink] (7.450,2.341) circle (1.1mm);
  \fill[ink] (6.959,3.610) circle (1.1mm);
  \fill[ink] (11.723,5.127) circle (1.1mm);
  \fill[ink] (8.647,2.375) circle (1.1mm);
  \fill[ink] (8.726,3.453) circle (1.1mm);
  \fill[ink] (1.172,4.401) circle (1.1mm);
  \fill[ink] (5.779,1.489) circle (1.1mm);
  \fill[ink] (7.328,5.355) circle (1.1mm);
  \fill[ink] (3.168,0.263) circle (1.1mm);
  \fill[ink] (3.692,3.998) circle (1.1mm);
  \fill[ink] (2.550,1.150) circle (1.1mm);
  \fill[ink] (10.706,3.896) circle (1.1mm);
  \fill[ink] (5.326,5.194) circle (1.1mm);
  \fill[ink] (3.993,3.929) circle (1.1mm);
  \fill[ink] (2.503,2.613) circle (1.1mm);
  \fill[ink] (10.411,2.353) circle (1.1mm);
  \fill[ink] (9.910,4.953) circle (1.1mm);
  \fill[ink] (3.411,1.147) circle (1.1mm);
  \fill[ink] (3.858,4.899) circle (1.1mm);
  \fill[ink] (11.592,2.734) circle (1.1mm);
  \fill[ink] (1.066,0.376) circle (1.1mm);
  \fill[ink] (10.325,0.432) circle (1.1mm);
  \fill[ink] (3.785,4.632) circle (1.1mm);
  \fill[ink] (0.422,0.961) circle (1.1mm);
  \fill[ink] (9.824,1.529) circle (1.1mm);
  \fill[ink] (1.834,0.463) circle (1.1mm);
  \fill[ink] (7.498,2.700) circle (1.1mm);
  \fill[ink] (7.508,3.868) circle (1.1mm);
  \fill[ink] (9.566,5.567) circle (1.1mm);
  \fill[ink] (8.140,1.316) circle (1.1mm);
  \fill[ink] (5.712,1.201) circle (1.1mm);
  \fill[ink] (8.484,1.203) circle (1.1mm);
  \fill[ink] (3.359,2.136) circle (1.1mm);
  \fill[ink] (7.328,4.435) circle (1.1mm);
  \fill[ink] (4.765,4.635) circle (1.1mm);
  \fill[ink] (10.712,0.688) circle (1.1mm);
  \fill[ink] (11.018,4.245) circle (1.1mm);
  \fill[ink] (1.707,2.740) circle (1.1mm);
  \fill[ink] (4.571,3.385) circle (1.1mm);
  \fill[ink] (10.400,4.662) circle (1.1mm);
  \fill[ink] (7.755,1.347) circle (1.1mm);
  \fill[ink] (8.574,4.783) circle (1.1mm);
  \fill[ink] (7.643,4.219) circle (1.1mm);
  \fill[ink] (2.674,5.240) circle (1.1mm);
  \fill[ink] (6.429,4.628) circle (1.1mm);
  \fill[ink] (1.160,2.669) circle (1.1mm);
  \fill[ink] (9.586,0.559) circle (1.1mm);
  \fill[ink] (4.086,4.612) circle (1.1mm);
  \fill[ink] (1.830,1.033) circle (1.1mm);
  \fill[ink] (6.192,4.252) circle (1.1mm);
  \fill[ink] (11.210,0.682) circle (1.1mm);
  \fill[ink] (3.566,4.282) circle (1.1mm);
  \fill[ink] (9.986,3.336) circle (1.1mm);
  \fill[ink] (3.816,2.335) circle (1.1mm);
  \fill[ink] (10.005,5.243) circle (1.1mm);
  \fill[ink] (2.616,4.964) circle (1.1mm);
  \fill[ink] (6.847,1.325) circle (1.1mm);
  \fill[ink] (7.221,0.355) circle (1.1mm);
  \fill[ink] (6.449,2.517) circle (1.1mm);
  \fill[ink] (11.300,5.371) circle (1.1mm);
  \fill[ink] (9.662,5.243) circle (1.1mm);
  \fill[ink] (5.728,1.976) circle (1.1mm);
  \fill[ink] (2.421,3.660) circle (1.1mm);
  \fill[ink] (10.933,0.925) circle (1.1mm);
  \fill[ink] (9.240,0.328) circle (1.1mm);
  \fill[ink] (8.170,2.004) circle (1.1mm);
  \fill[ink] (4.322,3.671) circle (1.1mm);
  \fill[ink] (1.417,4.293) circle (1.1mm);
  \fill[ink] (6.340,1.094) circle (1.1mm);
  \fill[ink] (7.606,3.166) circle (1.1mm);
  \fill[ink] (10.138,1.503) circle (1.1mm);
  \fill[ink] (10.671,1.969) circle (1.1mm);
  \fill[ink] (3.854,5.367) circle (1.1mm);
  \fill[ink] (2.730,5.791) circle (1.1mm);
  \fill[ink] (10.495,0.950) circle (1.1mm);
  \fill[ink] (3.210,0.743) circle (1.1mm);
  \fill[ink] (9.853,2.561) circle (1.1mm);
  \fill[ink] (9.363,0.906) circle (1.1mm);
  \fill[ink] (4.872,4.037) circle (1.1mm);
  \fill[ink] (8.116,5.304) circle (1.1mm);
  \fill[ink] (11.433,0.846) circle (1.1mm);
  \fill[ink] (2.392,0.595) circle (1.1mm);
  \fill[ink] (9.946,5.746) circle (1.1mm);
  \fill[ink] (0.919,5.528) circle (1.1mm);
  \fill[ink] (5.560,4.482) circle (1.1mm);
  \fill[ink] (3.991,2.815) circle (1.1mm);
  \fill[ink] (6.177,2.608) circle (1.1mm);
  \fill[ink] (8.332,4.928) circle (1.1mm);
  \fill[ink] (3.196,5.303) circle (1.1mm);
  \fill[ink] (8.835,4.557) circle (1.1mm);
  \fill[ink] (8.260,4.497) circle (1.1mm);
  \fill[ink] (8.583,0.595) circle (1.1mm);
  \fill[ink] (0.323,2.192) circle (1.1mm);
  \fill[ink] (7.927,2.092) circle (1.1mm);
  \fill[ink] (7.853,3.390) circle (1.1mm);
  \fill[ink] (11.569,0.328) circle (1.1mm);
  \fill[ink] (3.177,2.449) circle (1.1mm);
  \fill[ink] (0.785,1.294) circle (1.1mm);
  \fill[ink] (4.558,0.751) circle (1.1mm);
  \fill[ink] (11.419,3.381) circle (1.1mm);
  \fill[ink] (11.165,2.556) circle (1.1mm);
  \fill[ink] (6.310,3.548) circle (1.1mm);
  \fill[ink] (4.441,1.800) circle (1.1mm);
  \fill[ink] (3.634,0.278) circle (1.1mm);
  \fill[ink] (2.016,4.426) circle (1.1mm);
  \fill[ink] (4.723,5.226) circle (1.1mm);
  \fill[ink] (1.053,5.271) circle (1.1mm);
  \fill[ink] (11.489,1.565) circle (1.1mm);
  \fill[ink] (6.271,5.449) circle (1.1mm);
  \fill[ink] (8.584,2.823) circle (1.1mm);
  \fill[ink] (7.202,0.845) circle (1.1mm);
\end{tikzpicture}
\end{center}


(a) Count the dots in squares \textbf{B}, \textbf{C} and \textbf{F}. B: \wline[1.1cm] \quad C: \wline[1.1cm] \quad F: \wline[1.1cm]\\
(b) Average per square: \wline[1.8cm] \quad (c) Estimate for the whole field (average $\times$ 8): \wline[1.8cm]\\
(d) Now count every dot if you dare. Was your estimate close? \wline[1.8cm]
\end{puzzle}
\clearpage
